\documentclass[10pt,a4paper]{article}
\usepackage{amsmath,amssymb,amsthm,graphicx}
\usepackage{amsmath,amssymb,amsthm,hyperref,xfrac}
\providecommand{\eg}{e.g.\ }
\providecommand{\ie}{i.e.\ }
\usepackage[scr=rsfs,cal=euler]{mathalfa}
\multlinegap0pt
\hyphenation{accu-rate applied assump-tion assump-tions decom-pos-ing defined esti-mate esti-mates exten-ded hypothe-ses para-me-ter para-me-ters remain-der required}
\newenvironment{enumeratei}
{\bgroup\def\theenumi{\roman{enumi}}\def\theenumii{\arabic{enumii}}\begin{enumerate}}
{\end{enumerate}\egroup}
\newenvironment{enumeratea}
{\bgroup\def\theenumi{\alph{enumi}}\begin{enumerate}}
{\end{enumerate}\egroup}

\newcommand{\Psfrac}[2]{(\sfrac{#1}{#2})}
\newcommand{\psfrac}[2]{\sfrac{(#1)}{#2}}
\newcommand{\spfrac}[2]{\sfrac{#1}{(#2)}}
\newcommand{\pspfrac}[2]{\sfrac{(#1)}{(#2)}}
\newcommand{\shalf}{\sfrac{1}{2}}

\newcommand{\bnorm}[1]{\bigl\|#1\bigr\|}
\newcommand{\Bnorm}[1]{\Bigl\|#1\Bigr\|}
\newcommand{\babs}[1]{\bigl|#1\bigr|}
\newcommand{\Babs}[1]{\Bigl|#1\Bigr|}

\AtBeginDocument{%
\DeclareMathOperator{\sinn}{sin}\let\sin\sinn
\DeclareMathOperator{\coos}{cos}\let\cos\coos
}

%\usepackage{enumitem}
\usepackage{physics}
%\usepackage[foot]{amsaddr}
\usepackage{bm}
\usepackage{tikz}
\usepackage{mathtools}
\usepackage{scalerel,stackengine}
\stackMath
\usetikzlibrary{arrows.meta}
\usepackage{subcaption}
%\usepackage[style=numeric]{biblatex}
%\addbibresource{main.bib}
%\setlength{\bibitemsep}{0.3\baselineskip}

\newenvironment{mycases}
{\begin{dcases*}}
{\end{dcases*}}
\theoremstyle{plain}
\newtheorem{theorem}{Theorem}[section]
\newtheorem{lemma}[theorem]{Lemma}
\newtheorem{definition}[theorem]{Definition}
\newtheorem{proposition}[theorem]{Proposition}
\newtheorem{corollary}[theorem]{Corollary}
\newtheorem*{corollary*}{Corollary}
\theoremstyle{remark}
\newtheorem{remark}[theorem]{Remark}
\counterwithout{figure}{section}
\counterwithout{equation}{section}

\newcommand{\br}[1]{\left(#1\right)}
\newcommand{\bbr}[1]{\bigl(#1 \bigr)}
\newcommand{\Bbr}[1]{\Bigl(#1 \Bigr)}
\newcommand{\bbbr}[1]{\biggl(#1 \biggr)}
\newcommand{\tpd}{\frac{1}{(2\pi)^{d}}}
\newcommand\pig[1]{\scalerel*[2.35pt]{\big#1}{%
\ensurestackMath{\addstackgap[1.2pt]{\big#1}}}}
\newcommand\pigl[1]{\mathopen{\pig{#1}}}
\newcommand\pigr[1]{\mathclose{\pig{#1}}}
\newcommand{\ak}{\abs{k}}
\newcommand{\ap}{\abs{p}}
\newcommand{\aal}{\abs{\alpha}}
\newcommand{\sbbr}[1]{\bigl[#1\bigr]}
\newcommand{\sBbr}[1]{\Bigl[#1\Bigr]}
\newcommand{\sbbbr}[1]{\biggl[#1\biggr]}
\newcommand{\sbr}[1]{\left[#1\right]}
\newcommand{\half}{\frac{1}{2}}
\newcommand{\p}{\partial}
\newcommand{\jb}[1]{\langle#1\rangle}
\newcommand{\bset}[1]{\bigl\{#1\bigr\}}
\newcommand{\Bset}[1]{\Bigl\{#1\Bigr\}}
\newcommand{\set}[1]{\left\{#1\right\}}
\newcommand{\mytextfrac}[2]{{\textstyle\frac{#1}{#2}}}
\newcommand{\qaq}{\quad\text{and}\quad}
\newcommand{\para}{\parallel}
\newcommand{\ep}{\varepsilon}
\newcommand{\R}{\mathbb{R}}
\newcommand{\C}{\mathbb{C}}
\newcommand{\al}{\alpha}
\newcommand{\intRd}{\int_{\R^d}}
\newcommand{\intRdk}{\int_{\R^d_k}}
\newcommand{\intRdp}{\int_{\R^d_p}}
\newcommand{\lam}{\lambda}
\newcommand{\ceils}[1]{\lceil#1\rceil}
\newcommand{\Qin}{Q_{\mathrm{in}}}
\newcommand{\hin}{h_{\mathrm{in}}}
\newcommand\reallywidehat[1]{%
\savestack{\tmpbox}{\stretchto{%
\scaleto{%
\scalerel*[\widthof{\ensuremath{#1}}]{\kern.pt\mathchar"0362\kern.1pt}%
{\rule{0ex}{\textheight}}%WIDTH-LIMITED CIRCUMFLEX
}{\textheight}%
}{2.4ex}}%
\stackon[-6.9pt]{#1}{\tmpbox}%
}
\title{Uniform two-sided resonance-kernel integral bounds in dimensions at least three}
\author{Marnie Smith}
\date{doi:10.5802/jep.353; CC BY4.0}
\begin{document}
\maketitle
\noindent\textit{Editorial note.} Only the two original integral inequalities(one),(two) in the Control of resonances lemma are counted. The direct Schur operator corollary1571 is omitted from this counted unit rather than promoted to another task.\par\medskip
\noindent\textit{Editorial note.} Source1564 explicitly credits analogous estimates in Bedrossian-Masmoudi-Mouhot2018, Lemmas3.1 and3.2, and describes this proof as the extension to d>=3. Preserve that mathematical attribution and compare the exact generalized scope before uniqueness admission.\par\medskip
\noindent\textit{Editorial note.} The symbol x introduced at1584 denotes the auxiliary exponent1/(1-zeta), not a spatial integration variable. The implicit constants may depend on fixed d,m,beta and the chosen exponents, but not on T-star,t,s,k,ell.\par\medskip
\noindent\textit{Editorial note.} The changing-time integrals near1597/1627 use ordinary improper-integral estimates. No generated expansion of those estimates is being presented as author text.\par\medskip
\noindent\textit{Editorial scope.} One counted pair of Schur integral estimates. The immediately stated operator corollary1571, full nonlinear Hartree theorem and physical-space decay corollary are not separately counted.\par\medskip
\noindent\textit{Allowed standard prerequisites (editorial).}\begin{itemize}
\item Elementary weighted convolution estimates. Split Euclidean integrals into a neighborhood of each center and their complement; power tails are integrable when their combined exponent exceeds the dimension. Singular factors|x|\textasciicircum{}(-1/2) are locally integrable in the stated dimensions.. Check the one-dimensional parallel integrals and full-dimensional nonresonant integrals after the author changes of variables. No time-response kernel bound is supplied.
\item Tonelli, Cauchy-Schwarz and linear changes of variables. Integrate the positive kernel and rescale time along a nonzero parallel frequency; zero frequencies are handled separately.
\end{itemize}\medskip
\subsection{Quantum operators and norms} To formulate phase-mixing and scattering estimates in the quantum setting, quantum analogues of classical derivatives and weighted Sobolev norms are introduced. The standard and double Japanese brackets are defined by $\jb{x}:=(1+\abs{x}^2)^{1/2}$ and $\jb{x,y}:=(1+\abs{x}^2+\abs{y}^2)^{1/2}$, respectively. For a multi-index $\alpha=(\alpha_1,\dots,\alpha_d)\in\mathbb N^d$, write
$\abs{\alpha}=\alpha_1+\dots+\alpha_d$ and $D^{\alpha}=\p_1^{\alpha_1}\cdots\p_d^{\alpha_d}$.

Throughout, integrals such as $\int_x\,dx$ denote integration over $\R^d$ in the $x$ variable. The following conventions are adopted for the Fourier transform of functions $g=g(v)$ and $f=f(x,v)$:\vspace{-3pt}
\begin{align*}
\widehat g(p):=\int_v e^{-ip\cdot v}g(v)\,dv,\qquad \widehat f(k,p):=\int_x\int_ve^{-ik\cdot x}e^{-ip\cdot v}f(x,v)\,dvdx.
\end{align*}

\begin{lemma}[Control of resonances]\label{lemmaschurcontrolofresonances} Fix dimension $d\geq 3$ and let $T^*\geq0$, $m\geq\max\{1,(d-3/2)/2\}$ and $\beta\geq d+6$. The following integral inequalities hold:
\begin{align}\label{one}
\sup_{t\in[0,T^*],k\in\R^d}\int_0^t\int_{\ell\in\R^d}\frac{\jb{ s }\ak^{1/2}\abs{\ell}^{1/2}}{\jb{\ell}^m\jb{k-\ell,kt-\ell s }^{\beta}}d\ell d s &\lesssim1
\intertext{and}
\label{two}\sup_{s\in[0,T^*],\ell\in\R^d}\int_{s}^{T^*}\int_{k\in\R^d}\frac{\jb{ s }\ak^{1/2}\abs{\ell}^{1/2}}{\jb{\ell}^m\jb{k-\ell,kt-\ell s }^{\beta}}dk d t &\lesssim1.
\end{align}

\end{lemma}

\begin{proof}
To establish inequality \eqref{one}, fix $t\in[0,T^*]$ and $k\in\R^d$. If $k=0$, the integral vanishes, so assume $k\neq0$. The integral over early times is controlled as follows:
\begin{align*}
\int_0^1\hspace{-1mm}\int_\ell\!\frac{\jb{ s }\ak^{1/2}\abs{\ell}^{1/2}}{\jb{\ell}^m\jb{k-\ell,kt-\ell s }^{\beta}}d\ell d s \lesssim\int_\ell\!\frac{\ak^{1/2}\abs{\ell}^{1/2}}{\jb{\ell}^m\jb{k-\ell}^{\beta}}d\ell\lesssim\int_\ell\!\frac{\abs{\ell}^{1/2}}{\jb{\ell}^{m-{1/2}}\jb{k-\ell}^{\beta-\shalf}}d\ell\lesssim1,
\end{align*}where $\jb{k}\lesssim\jb{\ell}\jb{k-\ell}$ was used.

Assume $T^*>1$. Decompose $\ell\in\R^d$ into parallel and orthogonal components relative to $k$ as $\ell_\para:=(k\cdot\ell)k/\ak^2$ and $\ell_\perp:=\ell-\ell_\para$. Define $I_R\subset\R^d$ as a cylinder around the line spanned by $k$ with radius shrinking over time, and set $I_{NR}=I_R^c$, where
\begin{align*}
I_R=\bset{\ell:\abs{\ell_\perp}<\sfrac{\ak^b}{\jb{ s }^\zeta}}\qaq I_{NR}=\bset{\ell:\abs{\ell_\perp}\geq\sfrac{\ak^b}{\jb{ s }^\zeta}},
\end{align*}with $b:=(1-\zeta)/2$, where $\zeta$ is chosen to satisfy
\begin{align*}
\frac{2}{d-\sfrac{1}{2}}<\zeta<1\qaq x:=\frac{1}{1-\zeta}<\frac{11}{2}.
\end{align*}
Such a choice is possible since $(1-\zeta)^{-1}\to(2d-1)/(2d-5)\leq5$ as $\zeta\downarrow2/(d-1/2)$, for all $d\geq3$.

The integral over the cylinder $I_R$ satisfies, using $\abs{\ell}^{1/2}\lesssim\abs{\ell_\para}^{1/2}+\abs{\ell_\perp}^{1/2}$,
\begin{align*}
\mathscr{I}_R&:=\int_1^t\int_{I_R}\frac{\jb{ s }\ak^{1/2}\abs{\ell}^{1/2}}{\jb{\ell}^m\jb{k-\ell,kt-\ell s}^{\beta}}\,d\ell d s\\
&\lesssim\int_1^t\int_{I_R}\frac{\jb{ s }\ak^{1/2}}{\jb{\ell_\para}^m\jb{k-\ell_\para,kt-\ell_\para s }^\beta}\bbr{\abs{\ell_\para}^{1/2}+\abs{\ell_\perp}^{1/2}}\,d\ell d s.
\end{align*}
Applying the bound for $\abs{\ell_\perp}$ on $I_R$ and the volume of its cross-section gives
\begin{align*}
\mathscr I_R\lesssim\int_1^t\int_{-\infty}^\infty\frac{\jb{ s }\ak^{1/2}}{\jb{\ell_\para}^m\jb{k-\ell_\para,kt-\ell_\para s }^\beta}\frac{\ak^{(d-1)b}}{\jb{ s }^{(d-1)\zeta}}\Bbr{\abs{\ell_\para}^{1/2}+\frac{\ak^{\sfrac{b}{2}}}{\jb{ s }^{\sfrac{\zeta}{2}}}}\,d\ell_\para d s,
\end{align*}
leading to two contributions $\mathscr I_1$ and $\mathscr I_2$. For $\mathscr I_1$, the change of variable $ s \mapsto s /\abs{\ell_\para}$ and $(d-1)\zeta\geq1$ yield convergence under the conditions $\beta+m-2(d-1)b>3/2$ and $\beta,m\geq1/2+(d-1)b$. The second term $\mathscr I_2$ converges due to the choice of $\zeta$, ensuring that $1/\jb{ s }^{(d-1/2)\zeta-1}$ is integrable, for $\beta+m>2+2(d-1/2)b$ and $m,\beta\geq 1/2+(d-1/2)b$.

On $I_{NR}$, for $s\geq1$, it~holds that $\abs{kt-\ell s }\gtrsim s \abs{\ell_\perp}\gtrsim\ak^b\jb{ s }^{1-\zeta}$, so
\begin{align*}
\mathscr I_{NR}&:=\int_1^t\int_{I_{NR}}\frac{\jb{ s }\ak^{1/2}\abs{\ell}^{1/2}}{\jb{\ell}^m\jb{k-\ell,kt-\ell s}^{\beta}}\,d\ell d s \\
&\lesssim\int_1^t\int_{I_{NR}}\frac{\jb{s}\ak^{1/2}\abs{\ell}^{1/2}}{\jb{\ell}^m\jb{k-\ell,kt-\ell s}^{\beta-\spfrac{1}{2b}}\ak^{1/2}\jb{s}^{\pspfrac{1-\zeta}{2b}}}\,d\ell d s .
\end{align*}
Using $(1-\zeta)/(2b)-1=0$, and applying the change of variable $s\mapsto s/\abs{\ell}$, the integral $\mathscr I_{NR}$ converges for $\beta+m-1/(2b)-1/2>d$.

\pagebreak[2]
To establish inequality \eqref{two}, it~is observed that:
\begin{multline*}
\sup_{s\in[0,T^*],\ell\in\R^d}\int_{s}^{T^*}\int_{k\in\R^d}\frac{\jb{ s }\ak^{1/2}\abs{\ell}^{1/2}}{\jb{\ell}^m\jb{k-\ell,kt-\ell s }^{\beta}}\,dk d t \\
\leq \sup_{s\in[0,T^*],\ell\in\R^d}\int_{0}^{T^*}\int_{k\in\R^d}\frac{\jb{t }\ak^{1/2}\abs{\ell}^{1/2}}{\jb{\ell}^m\jb{k-\ell,kt-\ell s }^{\beta}}\,dk d t.
\end{multline*}
For fixed $s \in [0, T^*]$ and $\ell \in \mathbb{R}^d$, the integral over early times is estimated as
\begin{align*}
\int_{0}^{1}\int_{k}\frac{\jb{t }\ak^{1/2}\abs{\ell}^{1/2}}{\jb{\ell}^m\jb{k-\ell,kt-\ell s }^{\beta}}dk d t\lesssim\int_k\frac{\ak^{1/2}\abs{\ell}^{1/2}}{\jb{\ell}^m\jb{k-\ell}^{\beta}}\,dk\lesssim\int_k\frac{\ak^{1/2}}{\jb{k}^{\beta}}\,dk\lesssim1.
\end{align*}

Assume $T^* > 1$. If $\ell=0$, the integral vanishes, so assume $\ell\neq0$. Decompose $k\in\R^d$ into parallel and orthogonal components relative to $\ell$, defining $k_\para:=(\ell\cdot k)\ell/\abs{\ell}^2$ and $k_\perp:=k-k_\para$, and partition
\begin{align*}
J_R=\bset{k:\abs{k_\perp}<\sfrac{\abs{\ell}^b}{\jb{ t }^\zeta}}\qaq J_{NR}=\bset{k:\abs{k_\perp}\geq\sfrac{\abs{\ell}^b}{\jb{ t }^\zeta}},
\end{align*}with $\zeta$, $b$ and $x$ as in the proof of~\eqref{one}.

The integral over $J_R$ satisfies
\begin{align*}
\mathscr J_{R}&:=\int_1^{T^*}\int_{J_{R}} \frac{\jb{t}\abs{\ell}^{1/2}}{\jb{\ell}^m\jb{k_\para-\ell,k_\para t-\ell s}^\beta}\bbr{\abs{k_\para}^{1/2}+\abs{k_\perp}^{1/2}}\,dk dt\\
&\lesssim\int_1^{T^*}\int_{\R} \frac{\abs{\ell}^{1/2}}{\jb{\ell}^m\jb{k_\para-\ell,k_\para t-\ell s}^\beta}\frac{\abs{\ell}^{(d-1)b}}{\jb{t}^{(d-1)\zeta-1}}\Bbr{\abs{k_\para}^{1/2}+\frac{\abs{\ell}^{\sfrac{b}{2}}}{\jb{t}^{\sfrac{\zeta}{2}}}}\,dk_\para dt,
\end{align*}
leading to two contributions, $\mathscr J_1$ and $\mathscr J_2$. Following the change of variable $t\mapsto t/\abs{k_\para}$ in $\mathscr J_1$, the integrals converge under the conditions $m\geq(d-1/2)b+1/2$ and $(d-1/2)\zeta - 1 > 1$.

On $J_{NR}$, for $t\geq1$, it~holds that $\abs{kt-\ell s}\gtrsim t\abs{k_\perp}\gtrsim\abs{\ell}^b\jb{t}^{1-\zeta}$, so
\begin{align*}
\mathscr J_{NR}&:=\int_1^{T^*}\int_{J_{NR}}\frac{\jb{t}\abs{\ell}^{1/2}\ak^{1/2}}{\jb{\ell}^m\jb{k-\ell,kt-\ell s}^{\beta}}\,dkdt\\
&\lesssim\int_1^{T^*}\int_{J_{NR}}\frac{\abs{\ell}^{1/2}\ak^{1/2}}{\jb{\ell}^m\jb{k-\ell,kt-\ell s}^{\beta-x}\abs{\ell}^{bx}\jb{t}^{x(1-\zeta)-1}}\,dkdt.
\end{align*}Since $x(1-\zeta)=1$ and $bx=1/2$, the change of variable $t\mapsto t/\ak$ shows that the integral is uniformly bounded provided $\beta-x-1/2>d$. By the choice $x<11/2$, this condition follows from $\beta\geq d+6$. The remaining conditions obtained above follow from $\beta\geq d+6$ and $m\geq\max\{1,(d-3/2)/2\}$.
\end{proof}

\begin{thebibliography}{99}
\bibitem{BedrossianMasmoudiMouhot2018} J. Bedrossian, N. Masmoudi \& C. Mouhot - “Landau damping in finite regularity for unconfined systems with screened interactions”, Comm. Pure Appl. Math. 71 (2018) no. 3, p. 537-576
\end{thebibliography}
\end{document}
